\section{The Divine Motherhood: Mary, Mother of God}

\definitionaxis{The Virgin Mary is truly \emph{Theotokos}, Mother of God, because the one whom she conceived, bore, and brought forth according to his humanity is the eternal Son of God, one divine person in two natures.}{The Council of Ephesus received St.~Cyril of Alexandria's Second Letter to Nestorius and condemned the denial formulated in the first of Cyril's anathematisms (431); the Definition of Chalcedon located the title within its confession of the one and same Son (451).}{Mary neither causes the divine nature nor precedes the Trinity. ``Mother of God'' names a real relation to the divine person according to his human birth; it protects Christ's personal unity rather than divinizing his mother.}

The divine motherhood is architectonic for every other Marian doctrine. If the child born of Mary were a human person merely joined to the Word, calling her Mother of God would be false and the actions of Jesus could not be predicated of the one eternal Son. If his humanity were only apparent, her maternity, his suffering, and the saving obedience enacted in human nature would be unreal. The title stands at the intersection of Christ's true divinity, true humanity, and personal unity.

\subsection{The biblical reality before the technical title}

Scripture does not use the Greek noun \emph{Theotokos}; it gives the reality the word later protects. The subject of the birth narratives is the same subject confessed in the Gospel's highest Christology. Matthew identifies Mary's child as Jesus, Emmanuel, and the one who saves his people from their sins (Matt 1:18--25). Luke identifies him as holy, Son of God, Davidic king, and Lord (Luke 1:31--35, 43). Paul says that \emph{God sent his Son}, born of a woman (Gal 4:4): the Son does not begin to exist when he is born from her. John says that the Word who was with God and was God became flesh and dwelt among us (John 1:1--14). The one who assumes is the eternal Word; the humanity assumed from Mary is real and his own.

Luke 1:43 is especially compressed: Elizabeth, filled with the Holy Spirit, asks why ``the mother of my Lord'' has come to her. The text need not be forced to contain every technical distinction of fifth-century Christology. It does directly predicate maternity of Mary with respect to the one personally identified as Lord. The Church's later term does not replace this scriptural grammar; it excludes interpretations that dissolve it.

The infancy narratives also guard against a false conclusion. The Spirit's overshadowing and the Son's divine identity are God's work; Mary receives and consents. She contributes true human substance to the Son, so that he is genuinely ``from the seed of David according to the flesh'' (Rom 1:3), but she does not originate his divine sonship. Temporal maternity and eternal generation answer different questions: the Son is eternally begotten of the Father as God and temporally born of Mary as man.

\subsection{From anti-docetic confession to \emph{Theotokos}}

Early Christian witness defended the reality of this birth against docetism. St.~Ignatius of Antioch says that Jesus Christ was truly born from a virgin, truly nailed in the flesh, and truly raised (\emph{Smyrnaeans} 1--3); in \emph{Ephesians} 18--19 he joins Mary's virginity and childbirth to God's saving dispensation. The polemical center is Christ, but Mary's real maternity is indispensable evidence that the Son truly assumed flesh.

St.~Irenaeus develops the same realism through recapitulation. The Word takes flesh from Mary and renews Adam's race from within; Mary's obedient faith answers Eve's disobedience (\emph{Against Heresies} III.21.10; III.22.4). Irenaeus is not being cited as though he had written the Ephesian decree. He witnesses to its underlying logic: the Savior is truly the pre-existent Word and truly receives human life from the Virgin.

By the fourth century \emph{Theotokos} had become an established confession in prayer and theology. St.~Gregory Nazianzen makes the Christological stakes explicit: whoever does not accept holy Mary as \emph{Theotokos} is separated from the Godhead; if someone says Christ passed through Mary as through a channel rather than being formed in her both divinely and humanly, that person likewise destroys the Incarnation (Epistle 101 to Cledonius). Gregory's second clause is as important as the first. The title requires a real human derivation and gestation, not the use of Mary as an external conduit.

\subsection{Ephesus: the person born is the Word}

The controversy associated with Nestorius concerned the unity in which divine and human predicates belong to Christ. Antiochene concern to preserve the integrity of the natures was legitimate; language that made the man born of Mary a distinct personal subject alongside the Word was not. A preferred title such as ``Mother of Christ'' could be orthodox if ``Christ'' unambiguously named the one divine person, but it could also conceal a division. \emph{Theotokos} became the test because maternity terminates in a person: the one born of Mary is God the Word incarnate.

Cyril's Second Letter to Nestorius, formally received at Ephesus, explains that the Word did not first receive an ordinary man and then dwell in him. The Word personally united to himself flesh endowed with a rational soul, and thus underwent human birth. The holy Virgin is called Mother of God not because the nature of the Word or his divinity took origin from her, but because from her came the holy body with a rational soul to which the Word was hypostatically united. Cyril's First Anathema gives the concentrated consequence: if anyone does not confess Emmanuel truly God and therefore the holy Virgin \emph{Theotokos}, since she bore according to the flesh the Word made flesh, let that person be anathema.

Ephesus should not be reduced to a Marian slogan or a story of popular enthusiasm deciding doctrine. Its council acts, letters, contested procedures, and later reception require historical care. The dogmatic center is nevertheless clear in the received Cyrillic texts and the subsequent reunion. The Formula of Union of 433 confesses the Lord Jesus Christ as perfect God and perfect man, consubstantial with the Father in divinity and with us in humanity, and acknowledges the Virgin as Mother of God because of the union without confusion.

\subsection{Chalcedon: one and the same Son in two natures}

Chalcedon in 451 supplies the balanced grammar that prevents both division and confusion. Following the Fathers, it confesses ``one and the same Son,'' perfect in divinity and perfect in humanity, consubstantial with the Father according to divinity and with us according to humanity, begotten before the ages from the Father according to divinity and, in the last days, born for us and our salvation from Mary the Virgin, Mother of God, according to humanity. He is acknowledged in two natures without confusion, change, division, or separation.

The qualifying phrase ``according to humanity'' is not a retreat from \emph{Theotokos}. It is the precise reason the title is true without implying that Mary generates the Godhead. Natures answer \emph{what} Christ is; the person answers \emph{who} acts and is born. The human nature has no independent human personal subject beside the Word. Therefore the Word can truly be said to be born, suffer, and die according to his humanity, while the divine nature remains impassible. This is the \emph{communication of idioms}: predicates proper to either nature are said of the one person, without transferring the properties of one nature into the other.

\subsection{Aquinas: maternity terminates in the person}

St.~Thomas treats the question directly in \emph{Summa theologiae} III, q.35, a.4. To be conceived and born belongs to a person according to a nature. Since the human nature was assumed from the first instant of conception by the divine person, the one conceived and born from Mary is truly God. Mary is therefore truly Mother of God. Thomas uses an ordinary truth about motherhood to explain extraordinary Christology: a mother is mother of the person she bears, not merely of a detachable nature.

This article depends upon the surrounding treatment. ST III, q.2 defends the hypostatic union; q.31 explains Christ's true bodily derivation from Adam, Abraham, David, and Mary; q.32 distinguishes the Holy Spirit's efficient action from Mary's material contribution; and q.35 treats real nativity. The Spirit is not called Christ's father because divine conception does not establish the personal likeness and generative relation that fatherhood names. Mary is truly mother because Christ receives his human nature from her through conception and birth (ST III, q.32, a.3). Thomas thereby preserves the Father's eternal paternity, the Spirit's creative mission, and Mary's temporal maternity without rivalry.

In \emph{Compendium theologiae} I, cc.222--223, Thomas explicitly draws the anti-Nestorian consequence. If Christ were two persons, the one born of Mary would not be the eternal Son; but because the human and divine natures subsist in one person, the Virgin is correctly called Mother of God. The title is not devotional excess added after Christology. It is a compact test of the identity of Christ's personal subject.

\subsection{What the dogma reveals}

\begin{enumerate}
  \item \textbf{About Christ:} the Redeemer is one divine person, not a partnership of Word and man. His human obedience, suffering, death, and Resurrection are saving acts of the Son in his assumed humanity.
  \item \textbf{About salvation:} what is not assumed is not healed. The Son enters human genealogy, gestation, birth, family, law, work, suffering, and death. Salvation is neither escape from embodiment nor divine theater.
  \item \textbf{About Mary:} her dignity is relational and received. She is Mother of God because of her Son's identity, not because she is a goddess, a fourth person of the Trinity, or mother of the divine essence.
  \item \textbf{About the Church:} Mary is both mother and believer. Her bodily maternity is unique; her faith and obedient reception are exemplary for the Church, which receives Christ and bears witness to him (LG 53, 63--65).
\end{enumerate}

Pius XI's \emph{Lux veritatis} (1931), commemorating Ephesus, and the modern magisterium repeatedly return to this Christological center. Vatican II deliberately begins its Marian chapter within the mystery of Christ and the Church, confessing her as Mother of God and Mother of the Redeemer before considering her role (LG 52--53). Correct Marian devotion follows the same order: Christians venerate Mary with the singular honor traditionally called \emph{hyperdulia}, but adoration or \emph{latria} belongs to God alone (LG 66; CCC 971).

\subsection{Errors excluded}

It is false to say that the dogma teaches that Mary existed before God, generated the Trinity, or supplied Christ's divinity. It is inadequate to translate \emph{Theotokos} merely as ``Christ-bearer'' if that phrase evades the child's divine identity. It is equally inadequate to use the title while imagining a humanity only instrumentally inhabited by the Word. Finally, divine motherhood does not by itself prove every later Marian privilege. It is their Christological root and fitting ground; each further dogma still requires its own evidence and controlling act.
